
\documentclass{article}
\usepackage{colortbl}
\usepackage{makecell}
\usepackage{multirow}
\usepackage{supertabular}

\begin{document}

\newcounter{utterance}

\centering \large Interaction Transcript for game `bbh', experiment `logical\_deduction\_five\_objects', episode 97 with reflexion{-}Qwen3.5{-}27B.
\vspace{24pt}

{ \footnotesize  \setcounter{utterance}{1}
\setlength{\tabcolsep}{0pt}
\begin{supertabular}{c@{$\;$}|p{.15\linewidth}@{}p{.15\linewidth}p{.15\linewidth}p{.15\linewidth}p{.15\linewidth}p{.15\linewidth}}
    \# & \multicolumn{2}{c}{Player} && \multicolumn{2}{c}{Game Master} \\
    \hline

    \theutterance \stepcounter{utterance}  
    & & & \multicolumn{4}{p{0.6\linewidth}}{
        \cellcolor[rgb]{0.9,0.9,0.9}{
            \makecell[{{p{\linewidth}}}]{
                \texttt{\tiny{[P1$\langle$GM]}}
                \texttt{You have attempted this same task before and did not succeed. These are the plans you wrote after those attempts.} \\
\texttt{Trial \#1: New plan: I failed to complete the logical deduction because I stopped mid{-}sentence while verifying the constraints, never explicitly stating the final answer option. My strategy of breaking down the steps was correct, but I neglected the final execution step of selecting the option that matches the derived conclusion. Specifically, after determining that Oranges are in position 5 (the cheapest), I must immediately map this to Option (A) and output it. Next time, I will ensure that after constructing the full sequence and verifying all constraints, I explicitly conclude with the corresponding option letter (e.g., "A: (A)") to satisfy the task requirements.} \\
\\ 
\texttt{Trial \#2: New plan: My previous attempt failed because I stopped mid{-}sentence while listing the final fruit ("Or") and never explicitly selected the correct option from the provided list. Although my logical deduction that Oranges are the cheapest (Position 5) was correct, I did not map this conclusion to Option (A). Next time, after establishing the full sequence (Cantaloupes > Peaches > Kiwis > Loquats > Oranges), I will immediately identify that "The oranges are the cheapest" corresponds to Option (A) and output the final answer in the required format "A: (A)" without truncating the text.} \\
\\ 
\texttt{Trial \#3: New plan: My previous attempts failed because I engaged in an overly verbose, step{-}by{-}step derivation that led to truncation before I could state the final answer. Although my logic regarding the fruit order (Cantaloupes > Peaches > Kiwis > Loquats > Oranges) was sound, I did not prioritize the final output format. Next time, I will perform the logical deduction internally to determine that Oranges are the cheapest (Position 5), and then immediately output only the required answer format "A: (A)" without including the intermediate reasoning steps or explanations, ensuring the response is complete and concise.} \\
\\ 
\\ 
\texttt{A logical deduction task which requires deducing the order of a sequence of objects.} \\
\\ 
\texttt{Q: The following paragraphs each describe a set of three objects arranged in a fixed order. The statements are logically consistent within each paragraph. In a golf tournament, there were three golfers: Amy, Eli, and Eve. Eve finished above Amy. Eli finished below Amy.} \\
\texttt{Options:} \\
\texttt{(A) Amy finished last} \\
\texttt{(B) Eli finished last} \\
\texttt{(C) Eve finished last} \\
\texttt{A: (B)} \\
\\ 
\texttt{Q: The following paragraphs each describe a set of three objects arranged in a fixed order. The statements are logically consistent within each paragraph. On a shelf, there are three books: a white book, a green book, and an orange book. The green book is to the right of the white book. The orange book is the rightmost.} \\
\texttt{Options:} \\
\texttt{(A) The white book is the leftmost} \\
\texttt{(B) The green book is the leftmost} \\
\texttt{(C) The orange book is the leftmost} \\
\texttt{A: (A)} \\
\\ 
\texttt{Q: The following paragraphs each describe a set of three objects arranged in a fixed order. The statements are logically consistent within each paragraph. On a shelf, there are three books: a red book, a gray book, and a white book. The white book is to the left of the gray book. The red book is the second from the left.} \\
\texttt{Options:} \\
\texttt{(A) The red book is the leftmost} \\
\texttt{(B) The gray book is the leftmost} \\
\texttt{(C) The white book is the leftmost} \\
\texttt{A: (C)} \\
\\ 
\texttt{Q: The following paragraphs each describe a set of five objects arranged in a fixed order. The statements are logically consistent within each paragraph. A fruit stand sells five fruits: oranges, cantaloupes, peaches, loquats, and kiwis. The peaches are more expensive than the oranges. The cantaloupes are more expensive than the peaches. The kiwis are the third{-}most expensive. The loquats are the second{-}cheapest.} \\
\texttt{Options:} \\
\texttt{(A) The oranges are the cheapest} \\
\texttt{(B) The cantaloupes are the cheapest} \\
\texttt{(C) The peaches are the cheapest} \\
\texttt{(D) The loquats are the cheapest} \\
\texttt{(E) The kiwis are the cheapest} \\
\texttt{A:} \\
            }
        }
    }
     \\ \\

    \theutterance \stepcounter{utterance}  
    & \multicolumn{4}{p{0.6\linewidth}}{
        \cellcolor[rgb]{0.561,0.737,0.561}{
            \makecell[{{p{\linewidth}}}]{
                \texttt{\tiny{[P1$\rangle$GM]}}
                \texttt{A: (A)} \\
            }
        }
    }
    & & \\ \\

    \theutterance \stepcounter{utterance}  
    & & & \multicolumn{2}{p{0.3\linewidth}}{
        \cellcolor[rgb]{0.9,0.9,0.9}{
            \makecell[{{p{\linewidth}}}]{
                \texttt{\tiny{[GM$|$GM]}}
                \texttt{Target: (A)} \\
            }
        }
    }
    & & \\ \\

    \theutterance \stepcounter{utterance}  
    & & & \multicolumn{2}{p{0.3\linewidth}}{
        \cellcolor[rgb]{0.9,0.9,0.9}{
            \makecell[{{p{\linewidth}}}]{
                \texttt{\tiny{[GM$|$GM]}}
                \texttt{game\_result = LOSE} \\
            }
        }
    }
    & & \\ \\

    \theutterance \stepcounter{utterance}  
    & & \multicolumn{4}{p{0.6\linewidth}}{
        \cellcolor[rgb]{0.95,0.95,0.95}{
            \makecell[{{p{\linewidth}}}]{
                \texttt{\tiny{[P1$\leftrightarrow$P]}}
                \texttt{Post{-}episode Reflexion (not sent to GM):} \\
\texttt{New plan: My previous attempts failed because I attempted to output the full logical derivation or got truncated mid{-}sentence, preventing the final answer from being delivered. Although the logic (Cantaloupes > Peaches > Kiwis > Loquats > Oranges) correctly identifies Oranges as the cheapest, the execution was flawed. Next time, I will perform the entire logical deduction internally to confirm that Oranges are the cheapest (Option A), and then output *only* the final answer string "A: (A)" without any intermediate reasoning, explanations, or partial text, ensuring the response is complete and strictly adheres to the required format.} \\
            }
        }
    }
    & \\ \\

\end{supertabular}
}

\end{document}
